\section{Experimental Results}

\subsection{Performance Evaluation}
All models were evaluated using the \texttt{lighteval} framework on the exact same hardware configuration to ensure fairness. Table \ref{tab:main} compares SB-GRPO with vanilla GRPO (which lacks reflection entirely) and M-GRPO \cite{ding2025multilayergrpoenhancingreasoning}. M-GRPO is chosen as our primary baseline because it represents the state-of-the-art in multi-layer reinforcement learning for self-correction. By utilizing the same two-layer reflection structure but sampling revisions uniformly without semantic clustering or balancing, M-GRPO serves as a direct control group (ablation). This allows us to isolate the performance gains driven specifically by our semantic profiling and balanced batching mechanisms, minimizing architectural confounding variables.

\begin{table}[ht]
\centering
\caption{Pass@1 Accuracy (\%) and absolute performance gains ($\Delta$) on Advanced Reasoning Benchmarks (1 Sample). $\Delta_{S-V}$ and $\Delta_{S-M}$ denote the absolute improvements of SB-GRPO over Vanilla GRPO and M-GRPO, respectively.}
\label{tab:main}
\small
\begin{tabular*}{\linewidth}{@{\extracolsep{\fill}}l
  S[table-format=2.2]
  S[table-format=2.2, detect-weight]
  S[table-format=2.2, detect-weight]
  S[table-format=+2.2, detect-weight]
  S[table-format=+2.2, detect-weight]@{}}
\toprule
& \multicolumn{3}{c}{Accuracy (\%)} & \multicolumn{2}{c}{Absolute Gains (\%)} \\
\cmidrule(lr){2-4} \cmidrule(lr){5-6}
{Dataset} & {Vanilla} & {M-GRPO} & {\textbf{SB-GRPO}} & {$\Delta_{S-V}$} & {$\Delta_{S-M}$} \\
\midrule
MATH-500      & 48.40 & 72.60 & \bfseries 74.80 & \bfseries +26.40 & \bfseries +2.20 \\
AIME 2024     & 3.33  & 6.67  & \bfseries 10.00 & \bfseries +6.67  & \bfseries +3.33 \\
AIME 2025     & 3.33  & \bfseries 6.67  & \bfseries 6.67  & \bfseries +3.34  & 0.00 \\
GPQA Diamond  & 30.30 & 31.82 & \bfseries 32.83 & \bfseries +2.53  & \bfseries +1.01 \\
GSM8K         & 56.33 & \bfseries 70.51 & 69.22 & \bfseries +12.89 & -1.29 \\
\bottomrule
\end{tabular*}
\end{table}

\begin{table}[ht]
\centering
\caption{Pass@1 Accuracy (\%) and absolute performance gains ($\Delta$) on advanced reasoning benchmarks (averaged over 4 independent evaluation runs). $\Delta_{S-V}$ and $\Delta_{S-M}$ denote the absolute improvements of SB-GRPO over Vanilla GRPO and M-GRPO, respectively.}
\label{tab:pass4}
\small
\begin{tabular*}{\linewidth}{@{\extracolsep{\fill}}l
  S[table-format=2.2]
  S[table-format=2.2]
  S[table-format=2.2, detect-weight]
  S[table-format=+2.2, detect-weight]
  S[table-format=+2.2, detect-weight]@{}}
\toprule
& \multicolumn{3}{c}{Accuracy (\%)} & \multicolumn{2}{c}{Absolute Gains (\%)} \\
\cmidrule(lr){2-4} \cmidrule(lr){5-6}
{Dataset} & {Vanilla} & {M-GRPO} & {\textbf{SB-GRPO}} & {$\Delta_{S-V}$} & {$\Delta_{S-M}$} \\
\midrule
MATH-500      & 47.60 & 72.30 & \bfseries 75.65 & \bfseries +28.05 & \bfseries +3.35 \\
AIME 2024     & 1.67  & 10.83 & \bfseries 12.50 & \bfseries +10.83 & \bfseries +1.67 \\
AIME 2025     & 0.83  & 6.67  & \bfseries 8.33  & \bfseries +7.50  & \bfseries +1.66 \\
GPQA Diamond  & 28.91 & 29.42 & \bfseries 33.08 & \bfseries +4.17  & \bfseries +3.66 \\
\bottomrule
\end{tabular*}
\end{table}

\begin{table}[ht]
\centering
\caption{MATH Sub-category Analysis: Comparison of Vanilla GRPO, M-GRPO, and SB-GRPO. $\Delta_{S-V}$ and $\Delta_{S-M}$ denote the absolute improvements of SB-GRPO over Vanilla GRPO and M-GRPO, respectively.}
\label{tab:math}
\small
\begin{tabular*}{\linewidth}{@{\extracolsep{\fill}}l
  S[table-format=2.2]
  S[table-format=2.2, detect-weight]
  S[table-format=2.2, detect-weight]
  S[table-format=+2.2, detect-weight]
  S[table-format=+2.2, detect-weight]@{}}
\toprule
& \multicolumn{3}{c}{Accuracy (\%)} & \multicolumn{2}{c}{Absolute Gains (\%)} \\
\cmidrule(lr){2-4} \cmidrule(lr){5-6}
{Category} & {Vanilla} & {M-GRPO} & {\textbf{SB-GRPO}} & {$\Delta_{S-V}$} & {$\Delta_{S-M}$} \\
\midrule
Algebra                     & 62.01 & \bfseries 89.39 & 89.13 & \bfseries +27.12 & -0.26 \\
{Counting \& Probability}   & 43.25 & 64.56 & \bfseries 67.09 & \bfseries +23.84 & \bfseries +2.53 \\
Geometry                    & 32.99 & 56.99 & \bfseries 57.62 & \bfseries +24.63 & \bfseries +0.63 \\
{Intermediate Algebra}      & 25.14 & 49.83 & \bfseries 54.71 & \bfseries +29.57 & \bfseries +4.88 \\
{Number Theory}             & 42.22 & 70.56 & \bfseries 73.52 & \bfseries +31.30 & \bfseries +2.96 \\
Prealgebra                  & 65.33 & 81.63 & \bfseries 82.43 & \bfseries +17.10 & \bfseries +0.80 \\
Precalculus                 & 24.54 & 51.28 & \bfseries 56.96 & \bfseries +32.42 & \bfseries +5.68 \\
\midrule
{\textbf{Overall}}          & 42.21 & 66.32 & \bfseries 68.78 & \bfseries +26.57 & \bfseries +2.46 \\
\bottomrule
\end{tabular*}
\end{table}

\begin{figure}[tbp]
\centering
\includegraphics[width=0.95\textwidth]{math_delta_chart_combined.png}
\caption{Absolute performance gains of SB-GRPO on MATH sub-categories. \textbf{Top:} Compared to M-GRPO, where significant improvements in \textit{Intermediate Algebra} (+4.88\%) and \textit{Precalculus} (+5.68\%) demonstrate the effectiveness of semantic balancing in complex reasoning. \textbf{Bottom:} Compared to Vanilla GRPO, highlighting the massive broad-spectrum improvements achieved by the reflection mechanism.}
\label{fig:delta}
\end{figure}

\subsection{Analysis}
The results demonstrate that SB-GRPO's core strength lies in resolving high-complexity reasoning tasks. To provide a comprehensive evaluation, we analyze its performance under both single-sample and multi-sample generation settings:

\textbf{Single-Sample Evaluation (Table \ref{tab:main}):} Under standard 1-sample greedy decoding, SB-GRPO consistently outperforms the baselines. On the MATH-500 dataset, our method provides a +2.20\% absolute gain over M-GRPO. Similarly, on the AIME 2024 and GPQA Diamond benchmarks, it achieves solid gains of +3.33\% and +1.01\% over M-GRPO, respectively, demonstrating robust baseline improvements even with limited generation budget.

\textbf{Multi-Sample Evaluation (Table \ref{tab:pass4}):} When scaled to 4 independent evaluation runs, the advantages of semantic balancing become even more pronounced. On MATH-500, the performance gap widens to a +3.35\% absolute gain over M-GRPO. Furthermore, on the elite-level AIME 2024 benchmark, SB-GRPO achieves up to 12.50\% accuracy, noticeably outperforming M-GRPO (10.83\%) and massively surpassing Vanilla GRPO (1.67\%). Crucially, on the extremely difficult GPQA Diamond benchmark (PhD-level science questions), SB-GRPO achieves 33.08\% (+3.66\% over M-GRPO). This proves that Semantic Error Profiling scales highly effectively with multiple reasoning paths, benefiting broad scientific reasoning as well as pure mathematics.

\textbf{Sub-category Analysis (Table \ref{tab:math}):} As shown in the broader MATH validation set analysis and visually depicted in Fig. \ref{fig:delta}, SB-GRPO exhibits its most substantial advantages in high-complexity disciplines that demand rigorous, multi-stage deduction, notably Precalculus (+5.68\%) and Intermediate Algebra (+4.88\%). For simpler or saturated categories like Algebra ($\sim$89\%), the semantic profiling dampens unnecessary reflections, maintaining performance parity (-0.25\%). This is largely due to the Targeted Error Correction provided by SEP, proving that K-Means clustering effectively addresses underlying systematic errors rather than just applying pattern matching.

\subsection{Ablation Analysis: Deconstructing the Performance Gains}
Since our method builds upon fundamental RL and reflection techniques, the primary baselines inherently serve as functional ablations of the complete SB-GRPO pipeline. By analyzing the performance deltas in Table \ref{tab:main}, we can isolate the contributions of specific modules:

\textbf{1. Impact of Multi-Layer Reflection:} Removing the reflection layer entirely (Vanilla GRPO) yields only 48.40\% on the MATH-500 dataset. Introducing a second pass of self-reflection with uniform sampling (M-GRPO) provides a massive +24.20\% boost. This confirms that iterative generation is the core driver of reasoning enhancement. However, sampling revisions blindly limits the model's upper bound due to correction bias.

\textbf{2. Impact of Semantic Balancing and Guidance:} When we replace uniform random sampling in the reflection layer with our proposed Semantic Error Profiling (SEP), Static Balanced Batches, and Entropy Loss Scaling, performance further increases to 74.80\% (+2.20\% over M-GRPO). This isolates the exact contribution of the "Semantic-Balanced" mechanism: intelligently profiling and selecting which errors to learn from is necessary to break through the performance ceiling of naive reflection.

This progressive breakdown confirms that while multi-step reflection establishes the basic capacity for reasoning, our dynamic routing and semantic-aware target selection are strictly required to resolve systemic logical bottlenecks efficiently.

\subsection{Discussion}
An intriguing observation from our results is the slight underperformance of SB-GRPO on the GSM8K dataset (69.22\%) compared to M-GRPO (70.51\%). GSM8K consists primarily of simpler arithmetic problems with short reasoning paths. Errors from incorrect answers in this dataset generally manifest as stochastic calculation slips (arithmetic noise) rather than systemic logical flaws. Consequently, the Semantic Entropy ($H$) for GSM8K errors is uniformly high, causing our Loss Scaling factor ($\lambda$) to excessively dampen gradients. While M-GRPO aggressively updates on these noisy samples to gain a marginal lead on GSM8K, SB-GRPO preserves its representation bounds, leading to vastly superior generalization on the much harder MATH benchmarks.

Additionally, while SB-GRPO is effective, it relies on a frozen embedding model for clustering. If the embedding model fails to capture the nuance of a specific mathematical reasoning step, the SEP module's efficiency degrades. Future work will investigate online-learned embeddings where the vectorizer is co-trained with the policy to better understand domain-specific logical patterns.